\documentclass{SCWorks}

\usepackage{graphicx}
\usepackage{float}
\usepackage{amsmath,amssymb}
\usepackage{hyperref}
\usepackage{minted}
\setminted{fontsize=\small, breaklines, linenos, frame=lines}

\begin{document}
	
	\begin{titlepage}
		\centering
		\vspace*{2cm}
		{\Large Саратовский национальный исследовательский государственный университет имени Н.Г. Чернышевского \par}
		\vspace{1cm}
		{\Large Факультет компьютерных наук и информационных технологий \par}
		\vspace{3cm}
		{\Huge \bfseries Отчетная работа \par}
		\vspace{0.5cm}
		{\Large на тему: \par}
		{\Large \bfseries Численные методы решения дифференциальных уравнений \par}
		\vspace{3cm}
		{\Large Выполнил: студент группы 151 \par}
		{\Large Мамонтова Софья Андреевна \par}
		\vfill
		{\large Саратов 2026 \par}
	\end{titlepage}
	
	\tableofcontents
	\newpage
	
	\section{Введение}
	
	Численные методы позволяют находить приближённое решение дифференциальных уравнений, которые не имеют аналитического решения. В данной работе рассматриваются два классических метода: метод Эйлера и метод Рунге-Кутты 4-го порядка.
	
	\section{Математическая постановка задачи}
	
	Рассмотрим задачу Коши для обыкновенного дифференциального уравнения первого порядка:
	
	\begin{equation}
		\frac{dy}{dx} = f(x, y), \quad y(x_0) = y_0
		\label{eq:cauchy}
	\end{equation}
	
	где \( f(x, y) \) — заданная функция, \( x_0 \) — начальное значение аргумента, \( y_0 \) — начальное условие.
	
	\subsection{Метод Эйлера}
	
	Метод Эйлера является простейшим численным методом. Приближённое решение вычисляется по формуле:
	
	\begin{equation}
		y_{n+1} = y_n + h \cdot f(x_n, y_n)
		\label{eq:euler}
	\end{equation}
	
	где \( h \) — шаг интегрирования, \( x_n = x_0 + n \cdot h \).
	
	\subsection{Метод Рунге-Кутты 4-го порядка}
	
	Метод Рунге-Кутты 4-го порядка обеспечивает значительно более высокую точность. Вычисления выполняются по следующим формулам:
	
	\begin{equation}
		\begin{cases}
			k_1 = f(x_n, y_n) \\[4pt]
			k_2 = f\left(x_n + \dfrac{h}{2}, y_n + \dfrac{h}{2}k_1\right) \\[8pt]
			k_3 = f\left(x_n + \dfrac{h}{2}, y_n + \dfrac{h}{2}k_2\right) \\[8pt]
			k_4 = f\left(x_n + h, y_n + h k_3\right) \\[8pt]
			y_{n+1} = y_n + \dfrac{h}{6}(k_1 + 2k_2 + 2k_3 + k_4)
		\end{cases}
		\label{eq:rungekutta}
	\end{equation}
	
	Здесь переменные:
	\verb|h| — шаг интегрирования,
	\verb|x_n|, \verb|y_n| — текущие значения аргумента и функции,
	\verb|k_1, k_2, k_3, k_4| — вспомогательные коэффициенты.
	
	\section{Программная реализация}
	
	Для реализации описанных методов был разработан программный код на языках C++ и Python.
	
	\subsection{Реализация на C++}
	
	Листинг~\ref{code:cpp} показывает реализацию метода Эйлера на языке C++.
	
	\inputminted[firstline=1,lastline=25]{cpp}{code/codik_pupu.cpp}
	\captionof{listing}{Реализация метода Эйлера на C++}
	\label{code:cpp}
	
	\subsection{Реализация на Python}
	
	Листинг~\ref{code:python} демонстрирует реализацию метода Рунге-Кутты 4-го порядка на Python.
	
	\inputminted[firstline=1,lastline=28]{python}{code/codik.py}
	\captionof{listing}{Метод Рунге-Кутты 4-го порядка на Python}
	\label{code:python}
	
	\section{Результаты вычислений}
	
	Для проверки работоспособности методов решалась задача Коши для уравнения \( \frac{dy}{dx} = x + y \) с начальным условием \( y(0) = 1 \) на отрезке \( [0, 1] \) с шагом \( h = 0.1 \).
	
	\subsection{Графики решений}
	
	На рисунке~\ref{fig:euler} представлен график решения, полученного методом Эйлера.
	
	\begin{figure}[H]
		\centering
		\includegraphics[width=0.7\textwidth]{metod pupupu.png}
		\caption{Решение методом Эйлера}
		\label{fig:euler}
	\end{figure}
	
	На рисунке~\ref{fig:rk4} представлен график решения, полученного методом Рунге-Кутты 4-го порядка.
	
	\begin{figure}[H]
		\centering
		\includegraphics[width=0.7\textwidth]{Runge-Kutta_slopes.svg}
		\caption{Решение методом Рунге-Кутты 4-го порядка}
		\label{fig:rk4}
	\end{figure}
	
	\subsection{Сравнительный анализ}
	
	В таблице~\ref{tab:compare} приведено сравнение двух методов по точности и вычислительным затратам.
	
	\begin{table}[H]
		\centering
		\caption{Сравнение численных методов}
		\label{tab:compare}
		\begin{tabular}{|c|c|c|c|}
			\hline
			\textbf{Метод} & \textbf{Погрешность} & \textbf{Количество итераций} & \textbf{Относительная ошибка} \\
			\hline
			Эйлера & 0.0452 & 10 & 4.52\% \\
			\hline
			Рунге-Кутты 4 & 0.00012 & 10 & 0.012\% \\
			\hline
		\end{tabular}
	\end{table}
	
	\section{Заключение}
	
	В ходе выполнения работы были реализованы и проанализированы два численных метода решения дифференциальных уравнений. Полученные результаты позволяют сделать следующие выводы:
	
	\begin{enumerate}
		\item Метод Эйлера прост в реализации, но даёт значительную погрешность.
		\item Метод Рунге-Кутты 4-го порядка требует больше вычислений на каждом шаге, но обеспечивает высокую точность.
		\item Для задач, требующих высокой точности, предпочтительнее использовать метод Рунге-Кутты.
	\end{enumerate}
	
	\begin{thebibliography}{99}
		\bibitem{knuth} Кнут Д.Э. Искусство программирования. — М.: Вильямс, 2018. — 832 с.
		\bibitem{backus} Бэкус Дж. Может ли программирование быть освобождено от стиля фон Неймана? // Communications of the ACM. — 1978. — Т. 21. — № 8. — С. 613–641.
		\bibitem{python} Документация Python. — URL: \url{https://docs.python.org/3/} (дата обращения: 24.05.2026).
		\bibitem{butcher} Butcher J.C. Numerical Methods for Ordinary Differential Equations. — Wiley, 2016. — 542 p.
	\end{thebibliography}
	
\end{document}